\chapter{Apresentação}

O presente Projeto Pedagógico do Curso (PPC) de \ppccurso da Universidade Federal de Uberlândia foi elaborado para a implantação do currículo \ppcversao, com o propósito de consolidar uma formação em engenharia aplicada, crítica e socialmente engajada, alinhada às Diretrizes Curriculares Nacionais do Curso de Graduação em Engenharia~\cite{cne_ces_2_2019,cne_ces_1_2021} e às orientações institucionais da UFU~\cite{ufu_guia_ppc_2021}.

Este PPC responde à crescente demanda por profissionais capazes de projetar, desenvolver, integrar e gerir sistemas automatizados industriais em um cenário de transformação digital da indústria — a chamada Indústria 4.0 —, marcado pela convergência entre eletrônica, controle, redes industriais, robótica e inteligência artificial aplicada. O curso foi construído de forma coletiva pelo Núcleo Docente Estruturante (NDE), pelo Colegiado do Curso e por representantes da comunidade acadêmica, com a convicção de que a formação de um(a) \ppctitulacao deve ser tecnicamente rigorosa, profissionalizante e orientada por competências diretamente ligadas ao exercício profissional.

A base regulatória deste curso são as Diretrizes Curriculares Nacionais do Curso de Graduação em Engenharia~\cite{cne_ces_2_2019,cne_ces_1_2021}; o parágrafo herdado do Curso Superior de Tecnologia, que citava o CNCST, o CINE Brasil 0714A01 e a Resolução CNE/CP nº 2/2024, foi removido por não se aplicar a um Bacharelado em Engenharia (\autoref{chp:identificacao}). \textcolor{red}{[CONFIRMAR JUNTO AO NDE/COLEGIADO A DENOMINAÇÃO REGULATÓRIA ESPECÍFICA DESTE CURSO — NÃO HÁ, ATÉ A ELABORAÇÃO DESTE PPC, PORTARIA MEC OU RESOLUÇÃO CNE QUE INSTITUA ESPECIFICAMENTE A DENOMINAÇÃO ``ENGENHARIA DE AUTOMAÇÃO, ROBÓTICA E INTELIGÊNCIA ARTIFICIAL''; A DENOMINAÇÃO PRECEDENTE MAIS PRÓXIMA NO SISTEMA CONFEA/CREA E NO PRÓPRIO MEC É ``ENGENHARIA DE CONTROLE E AUTOMAÇÃO'' (INSTITUÍDA PELA PORTARIA MEC Nº 1.694/1994)]}

Este PPC reflete um currículo estruturado em núcleos de formação básica, tecnológica e profissional, humanística, de integração curricular/extensão e de formação complementar, combinando disciplinas de fundamentos científicos, tecnologias centrais de automação industrial (eletricidade, eletrônica, instrumentação, controle, CLP/SCADA, redes industriais, acionamentos, hidráulica e pneumática) e um eixo diferenciador em robótica, sistemas inteligentes e inteligência artificial aplicada à automação. Nos períodos específicos deste curso, os componentes curriculares optativos são organizados em Ênfases Formativas — Máquinas Inteligentes, Acionamentos e Proteção Industrial (MIAPI); Robótica Autônoma e Sistemas de Controle (RASC); e Sistemas Embarcados, IoT e Computação Inteligente (SEICI) —, permitindo ao(à) estudante aprofundar-se em, no mínimo, duas dessas ênfases (\autoref{chp:modulos_optativos}). Mais do que uma exigência normativa, este documento representa a carta de intenções de um curso que se projeta para atender à demanda do setor produtivo por profissionais de formação em engenharia aplicada, sem perder de vista o rigor técnico-científico e a responsabilidade social de sua atuação.

O \ppccurso compõe uma proposta de Área Básica de Ingresso (ABI) com o Curso Superior de Tecnologia em Automação Industrial: os cinco primeiros períodos das matrizes são comuns e a proposta prevê a escolha da trajetória após essa etapa. A organização preserva a duração e a identidade próprias de cada curso — o \ppccurso integraliza-se em \ppctempominimo, com cinco períodos específicos. O ingresso, a escolha e eventual continuidade em uma segunda graduação dependem de ato institucional específico, conforme o \autoref{chp:abi}.
